\section{Two adjacent real-rooted rows that fail to interlace}
The size-40 row in Theorem~\ref{single} fails real-rootedness. The next example isolates a failure of relative root order. It compares two rows whose individual roots are all real and simple.

\begin{theorem}\label{two-cycle}
Use the same deterministic length set at both sizes,
\[
 S_\ast=\{44,46,48,50,52,54,76,78,80,82,84,86\}\ \cup\ T_{87}.
\]
Equivalently, for every positive integer length $j$ its weight is one if $j\in S_\ast$ and zero otherwise. Every contributing permutation at size $129$ or $130$ has at most two cycles. The rows $F_{129}(S_\ast)$ and $F_{130}(S_\ast)$ each have one simple zero at zero, and respectively $127$ and $128$ simple negative zeros. They fail weak interlacing.
\end{theorem}

\begin{proof}
Every retained length is at least $44$, and $3\cdot44>130$. At size $129$, two short lengths would have even sum, and a length from $T_{87}$ and any retained length sum to at least $131$. Thus only the full $129$-cycle contributes. For size $130$, the two-cycle types are precisely $(a,130-a)$ for
\[
 \mathcal A=\{44,46,48,50,52,54\}.
\]
Their lengths differ, so choosing the support of the shorter cycle gives exactly $\binom{130}{a}$ choices for the two supports. The excedance counts on the two cycles add. By the single-cycle formula recorded before \eqref{rec}, the quotients $P=F_{129}(S_\ast)/t$ and $Q=F_{130}(S_\ast)/t$ are
\begin{equation}\label{two-cycle-rows}
 P=A_{128},\qquad
 Q=A_{129}+t\sum_{a\in\mathcal A}\binom{130}{a}A_{a-1}A_{129-a}.
\end{equation}
Their degrees are $127$ and $128$. Their leading and constant coefficients are one, and both are reciprocal by inversion.

To verify individual splitting, the attached exact certificate isolates $63$ roots of $P$ and $64$ of $Q$ in $(-1,0)$ using disjoint open dyadic brackets within each row. At each bracket the endpoint signs are opposite and nonzero. The odd-degree reciprocal polynomial $P$ has $P(-1)=0$, and exact evaluation gives $Q(-1)\ne0$. Reciprocals give equally many roots below $-1$, so the two degrees exhaust the root counts. All the negative roots are simple, and $P(0)=Q(0)=1$ makes the extra zero simple in both full rows.

For a direct obstruction to interlacing use the rational point
\[
 t_\ast=-1/9800000000000000 .
\]
Exact evaluation in \eqref{two-cycle-rows} gives the four signs
\[
 P(t_\ast)<0,\qquad P'(t_\ast)>0,\qquad Q(t_\ast)>0,\qquad Q'(t_\ast)>0 .
\]
For reproducibility, Table~\ref{two-cycle-signs} gives open rational enclosures. An integer method is to take $D=-1/t_\ast$, and, for $f=\sum_{i=0}^d f_i t^i$, compute
\[
 D^d f(-1/D)=\sum_{i=0}^d(-1)^i f_i D^{d-i},\qquad
 D^{d-1}f'(-1/D)=\sum_{i=1}^d(-1)^{i-1}i f_i D^{d-i}.
\]
The Eulerian coefficients in \eqref{two-cycle-rows} are generated by the recurrence following \eqref{rec}. The supplementary verifier also recomputes them by alternating binomial sums.

If the full rows weakly interlaced, deleting their common largest zero would give weak interlacing of $P$ and $Q$. At a root $\alpha$ of the monic $P$, the order of the roots would then give $Q(\alpha)/P'(\alpha)\le0$. Indeed, the signs of these two values are opposite whenever $Q(\alpha)\ne0$. Partial fractions give
\[
 (Q/P)'(t)=1-\sum_{P(\alpha)=0}\frac{Q(\alpha)/P'(\alpha)}{(t-\alpha)^2}
\]
away from the roots of $P$. Thus $Q'P-QP'$ would be positive at $t_\ast$. The four signs above make it strictly negative there, a contradiction.
\end{proof}

\begin{table}[ht]
\centering
\begin{tabular}{@{}cccc@{}}
\hline
$P(t_\ast)/10^{26}$ & $P'(t_\ast)/10^{45}$ &
$Q(t_\ast)/10^{26}$ & $Q'(t_\ast)/10^{45}$\\ \hline
$(-3,-2)$ & $(1,2)$ & $(1,2)$ & $(4,5)$\\ \hline
\end{tabular}
\caption{Open rational enclosures for the four exact sign tests.}\label{two-cycle-signs}
\end{table}

The supplementary dyadic computation also refines the brackets until the two rows' brackets are disjoint from one another. In $(-1,0)$, reading from zero toward $-1$, the $P,Q$ tags agree with the alternating list $Q,P,Q,P,\ldots,Q$ except that positions five and six exchange tags. Reciprocity reflects the exchange on the other side of $-1$. The formal proof uses the four rational signs for the pairwise obstruction and separate full-axis negative-root certificates for the two rows. The exact exchange of ranks is verified by the separate integer computation.

Theorem~\ref{two-cycle} concerns an independent enlargement of Alexandersson's hard minimum-length question. We use $87$ as the tail parameter, with twelve retained lengths below it. The least retained length is $44$, and the unit tail can also be written as $T_{86}$ because $86$ is retained. Taking the hard cutoff at $87$ would leave only the two full-cycle rows, since $130<2\cdot87$. Taking the hard cutoff at $44$ would also admit all intervening lengths. The hard-cutoff two-cycle theorem of ProgramComputer applies to the latter pair, since its upper size is at most $3\cdot44-1$ \cite{pc-two}. The selective profile $S_\ast$ gives a different polynomial even with at most two cycles at both sizes.

At fixed size eight, the rows with hard cutoffs four and five already fail to interlace \cite{earlier}. In Theorem~\ref{two-cycle} the length rule stays fixed as the size increases.
